\documentclass[12pt]{article}
\usepackage{amsmath}
\usepackage{amsfonts}
\usepackage{amssymb}
\usepackage{geometry}
\usepackage{tikz}
\geometry{a4paper, margin=1in}
\pagenumbering{gobble}

\title{02YMECH - Homework 3}
\date{Assigned for the week of Oct 6, 2025}
\author{}
\begin{document}
\maketitle

\section*{Questions}

\begin{enumerate}
    \item A child of mass $m_c$ is sitting in a boat of mass $m_b$, which is initially moving with velocity $\vec{v}_0$ on calm water. The child throws a package of mass $m_p$ with velocity $\vec{v}_p$ at an angle $\theta$ above the horizontal, measured relative to the water (i.e., an inertial frame). Assuming that no external horizontal forces act on the system and neglecting any water resistance, determine the velocity $\vec{v}_b$ of the boat (including both $x$ and $y$ components) immediately after the throw.
    \item A particle of mass $m$ moves along the $x$-axis under the influence of a position-dependent force given by
\[
F(x) = - \frac{k x}{(x^2 + a^2)^{3/2}},
\]
where $k>0$ and $a>0$ are constants. Initially, the particle is at rest at $x = x_0 > 0$. 

\begin{enumerate}
    \item Calculate the work done by the force as the particle moves from $x = x_0$ to a general position $x$. 
    \item Using the work-energy theorem, find the velocity of the particle as a function of $x$. 
    \item Determine the maximum speed of the particle. 
\end{enumerate}
    \item A solid disk of mass $M$ and radius $R$ is free to rotate about its central axis (perpendicular to its plane) and is initially at rest. A small particle of mass $m$ is moving with speed $v_0$ in a direction tangential to the edge of the disk and lands at the rim of the disk, sticking to it (i.e., the particle’s velocity is perpendicular to the radius vector connecting the disk’s center to the point of contact). The collision is perfectly inelastic. The moment of inertia of a solid disk about its central axis is $I = \frac{1}{2} M R^2$.
    \begin{enumerate}
        \item Determine the angular velocity $\omega$ of the disk+particle system immediately after the particle sticks.
        \item Determine the rotational kinetic energy $T_\text{rot}$ of the disk+particle system after the collision.
    \end{enumerate}
    \item A particle of mass $m$ moves along the $x$-axis under the influence of a conservative force derived from the potential
\[
U(x) = U_0 \left( \frac{x^4}{a^4} - 2 \frac{x^2}{a^2} \right),
\]
where $U_0 > 0$ and $a$ are constants.

\begin{enumerate}
    \item Determine the force acting on the particle as a function of $x$.
    \item Find the positions of stable and unstable equilibrium.
    \item If the particle is released from rest at $x = 0$, determine its speed when it reaches $x = a$.
    \item Compute the work done by the force as the particle moves from $x = 0$ to $x = a$.
\end{enumerate}

\end{enumerate}
\end{document}
